% Page 056

\seite{56}

\abschnitt{Schuljahr 1919–20}
\pstart
Entlassen wurden vierzig 2 Knaben, Ostern 1 Mädchen. Alle widmen sich\ob
der Landwirtschaft. Aufgenommen wurden 8 Kinder, 6 Mädchen, 2\ob
Knaben. Die Schülerzahl beträgt so jetzt 44 Schüler, davon 12 Knaben,\ob
32 Mädchen.

\pend
\abschnitt{Kriegsfolge 5}
\pstart

Der Krieg ist beendet, das Dorf hat 7 Gefallene zu beklagen. Die tapferen\ob
Helden werden wir nicht vergessen. Die anderen Krieger sind in die Heimat\ob
eingetroffen, d.\ h. ein Jüngling von 23 Jahren als Krüppel, ein Bein hat er\ob
verloren, 4 aus dem Dorfe waren in Gefangenschaft, einer ist in englischer\ob
zurück, 3 andere sind in französischer. Alle hoffen sie auf bald in der\ob
Heimat anlangen zu können.

\pend
\abschnitt{Ernte}
\pstart

Die Gesamternte lieferte stellenweise einen geringen Ertrag. Die Getreideernte\ob
war verhältnismäßig galt gegen letztes Jahr mehr Kartoffelfrüchte viel\ob
eingebracht. Die Preise im allgemeinen seigen fortwährlig, sowohl Vieh\ohb
preise als Butterpreise.

\pend
\abschnitt{Schuljahr 1920–21}
\pstart

Vorzeitig wurde Ostern Herbst 1919 ein Schüler entlassen, Ostern\ob
1920 1 Mädchen und 3 Knaben, sämtliche widmen sich der Land\ohb
wirtschaft. Aufgenommen wurden an Kinder: 4 Mädchen, 3 Kna\ohb
ben. Die Schülerzahl beträgt 24, 29 Mädchen, 15 Knaben.

\pend
\abschnitt{Ernte}
\pstart

In der hiesigen Gemeinde fiel allgemein befriedigend aus, so\ohb
wohl das Getreide, wie auch die\unsicher{} Kartoffelernte besser, als man erwar\ohb
tet hatte. Die Witterung war eine besonders günstige, namentlich\ob
auch für Kartoffeln. Die Preise sind für Getreide und Vieh sehr\ob
stark gefallen, es ist schlecht, infolge der Einführung Klauenseuche-Grenz\ohb
sperre sollte im Herbste große Märkte ausfallen, auch sollten von den\ob
besetzten Gebieten Lebensmittelzufuhr sind noch wie nur immer sehr schwach.
\pend
